% ============================================================
\chapter{GARCH Models and Conditional Volatility}
\label{ch:garch}
% ============================================================

\section{Motivating example}
Financial returns are approximately uncorrelated, but their squares $r_t^2$ exhibit strong autocorrelation: periods of high volatility are clustered (\textbf{volatility clustering}\index{volatility clustering}). ARMA models do not capture this phenomenon.

\section{Stylised facts of financial returns}\index{stylised facts}

\begin{enumerate}
\item Returns are approximately uncorrelated: $\rho_{r}(h)\approx 0$.
\item Squared returns are correlated: $\rho_{r^2}(h)>0$ (volatility persistence).
\item Heavy tails (leptokurtosis): kurtosis $>3$.
\item Volatility asymmetry (leverage effect): declines increase volatility more than rises.
\end{enumerate}

\section{ARCH model}\index{ARCH}

\begin{definition}[ARCH$(q)$]
\[
r_t = \sigma_t\,z_t, \quad z_t\iid(0,1), \quad \sigma_t^2 = \omega + \sum_{i=1}^q\alpha_i r_{t-i}^2,
\]
where $\omega>0$, $\alpha_i\geq 0$, and $\sum\alpha_i<1$ for stationarity.
\end{definition}

\begin{theorem}[Properties of an ARCH$(q)$]
\begin{enumerate}
\item $\E[r_t]=0$, $\Var(r_t)=\omega/(1-\sum\alpha_i)$.
\item $(r_t)$ is a white noise (uncorrelated).
\item $(r_t^2)$ follows an AR$(q)$: $r_t^2 = \omega+\sum\alpha_i r_{t-i}^2+v_t$ where $v_t=r_t^2-\sigma_t^2$ is a white noise.
\item Kurtosis $>3$ (heavy tails even if $z_t$ is Gaussian).
\end{enumerate}
\end{theorem}

\section{GARCH model}\index{GARCH}

\begin{definition}[GARCH$(p,q)$]
\[
r_t = \sigma_t z_t, \quad \sigma_t^2 = \omega + \sum_{i=1}^q\alpha_i r_{t-i}^2 + \sum_{j=1}^p\beta_j\sigma_{t-j}^2.
\]
Stationarity conditions: $\omega>0$, $\alpha_i,\beta_j\geq 0$, $\sum_{i=1}^q\alpha_i+\sum_{j=1}^p\beta_j<1$.
\end{definition}

\begin{theorem}[Unconditional variance of GARCH$(1,1)$]
\[
\Var(r_t) = \frac{\omega}{1-\alpha_1-\beta_1}.
\]
The \textbf{persistence} of volatility is measured by $\alpha_1+\beta_1$ (close to 1 in finance).
\end{theorem}

\begin{intuition}
The GARCH$(1,1)$ is the most widely used volatility model in practice. The one-step-ahead conditional variance forecast is:
\[
\hat{\sigma}_{t+1}^2 = \omega + \alpha_1 r_t^2 + \beta_1\sigma_t^2.
\]
It is an exponentially weighted moving average of past squared values.
\end{intuition}

\section{IGARCH}\index{IGARCH}

\begin{definition}
If $\alpha_1+\beta_1=1$ (integrated GARCH), volatility shocks persist indefinitely. The process is not second-order stationary but remains \textbf{strictly stationary}.
\end{definition}

\section{Extensions}

\subsection{EGARCH}\index{EGARCH}
\begin{definition}
Nelson's EGARCH model:
\[
\ln\sigma_t^2 = \omega + \alpha(|z_{t-1}|-\E[|z_{t-1}|]) + \gamma z_{t-1} + \beta\ln\sigma_{t-1}^2.
\]
The term $\gamma z_{t-1}$ captures the \textbf{leverage effect}. No positivity constraint on the parameters is needed.
\end{definition}

\subsection{GJR-GARCH}\index{GJR-GARCH}
\begin{definition}
\[
\sigma_t^2 = \omega + (\alpha+\gamma\mathbf{1}_{r_{t-1}<0})r_{t-1}^2 + \beta\sigma_{t-1}^2.
\]
$\gamma>0$ implies that negative returns increase volatility more.
\end{definition}

\section{Estimation}\index{estimation!GARCH}

\begin{definition}[Quasi-maximum likelihood]
The Gaussian conditional log-likelihood:
\[
\ell(\bm\theta) = -\frac{1}{2}\sum_{t=1}^n\left(\ln\sigma_t^2 + \frac{r_t^2}{\sigma_t^2}\right).
\]
Even if $z_t$ is not Gaussian, the QMLE is consistent under certain conditions.
\end{definition}

\section{Tests for ARCH effects}\index{test!ARCH}

\begin{definition}[Engle's ARCH-LM test]
Regress $r_t^2$ on $r_{t-1}^2,\ldots,r_{t-q}^2$ and test for joint significance via $nR^2\sim\chi^2(q)$.
\end{definition}

\section{Implementation}

\begin{codeblock}[title=\textbf{Python}]
\begin{minted}{python}
import numpy as np
import pandas as pd
import matplotlib.pyplot as plt
from arch import arch_model

# Simuler un GARCH(1,1)
np.random.seed(42)
n = 1000
omega, alpha, beta = 0.05, 0.10, 0.85
sigma2 = np.zeros(n)
r = np.zeros(n)
sigma2[0] = omega / (1 - alpha - beta)
for t in range(1, n):
    sigma2[t] = omega + alpha * r[t-1]**2 + beta * sigma2[t-1]
    r[t] = np.sqrt(sigma2[t]) * np.random.normal()

fig, axes = plt.subplots(3, 1, figsize=(12, 8))
axes[0].plot(r, lw=0.5)
axes[0].set_title('Rendements simulés (GARCH(1,1))')
axes[1].plot(r**2, lw=0.5, color='orange')
axes[1].set_title('$r_t^2$ (volatilité réalisée)')
axes[2].plot(np.sqrt(sigma2), lw=0.5, color='red')
axes[2].set_title('$\\sigma_t$ (volatilité conditionnelle)')
plt.tight_layout()
plt.savefig('ch06_garch.pdf')
plt.show()

# Ajuster un GARCH(1,1)
am = arch_model(r, vol='Garch', p=1, q=1, dist='normal')
res = am.fit(disp='off')
print(res.summary())

# Prévision de volatilité
forecasts = res.forecast(horizon=10)
print("Variance prévisionnelle:")
print(forecasts.variance.iloc[-1])
\end{minted}
\end{codeblock}

\section{Exercises}

\begin{exercise}[$\star$]
Show that an ARCH(1) with $r_t=\sigma_t z_t$, $\sigma_t^2=\omega+\alpha r_{t-1}^2$ implies $\E[r_t^4]=3\omega^2(1+\alpha)/(1-3\alpha^2)$ (if $3\alpha^2<1$). Deduce the kurtosis.
\end{exercise}

\begin{exercise}[$\star\star$]
Show that the GARCH$(1,1)$ can be written as an ARCH$(\infty)$.
\end{exercise}

\begin{exercise}[$\star\star$]
Prove that the process $r_t^2$ of a GARCH$(1,1)$ follows an ARMA$(1,1)$.
\end{exercise}

\begin{exercise}[$\star\star\star$ -- Project]
Download the daily returns of the S\&P~500 over 10 years. Estimate a GARCH$(1,1)$, an EGARCH$(1,1)$, and a GJR-GARCH$(1,1)$. Compare by AIC and analyse the leverage effect.
\end{exercise}

\begin{keyformulas}
\begin{align*}
\text{GARCH}(1,1) &: \sigma_t^2 = \omega+\alpha r_{t-1}^2+\beta\sigma_{t-1}^2, \quad \alpha+\beta<1 \\
\Var(r_t) &= \omega/(1-\alpha-\beta) \\
\text{EGARCH} &: \ln\sigma_t^2 = \omega+\alpha(|z_{t-1}|-\E|z_{t-1}|)+\gamma z_{t-1}+\beta\ln\sigma_{t-1}^2
\end{align*}
\end{keyformulas}
